\chapter{PENUTUP}
\label{sec:chap5_penutup}

\section{Kesimpulan}

Berdasarkan hasil perancangan, implementasi, dan pengujian yang telah dilakukan, dapat diperoleh beberapa kesimpulan sebagai berikut:

\begin{enumerate}
	
	\item Platform monitoring CCTV berbasis ESP32-CAM berhasil dirancang dan diimplementasikan sehingga dapat digunakan untuk melakukan pemantauan perangkat melalui website dan aplikasi mobile secara real-time. Sistem mampu menampilkan informasi kamera, data citra, status perangkat, serta aktivitas yang terdeteksi pada area pengawasan.
	
	\item Sistem komunikasi dan penyimpanan data berhasil diimplementasikan menggunakan MQTT, WebSocket, MySQL, dan MinIO yang dijalankan pada server Intel NUC. Penggunaan MQTT memungkinkan proses komunikasi data antara perangkat dan server berjalan dengan baik, sedangkan WebSocket mampu menyediakan pembaruan data secara real-time tanpa mekanisme polling sehingga meningkatkan efisiensi komunikasi dan pengalaman pengguna.
	
	\item Fitur motion detection, monitoring telemetri perangkat, dan notifikasi mobile berhasil diimplementasikan pada sistem. Motion detection mampu mendeteksi aktivitas pada area pengawasan dan menghasilkan event yang digunakan untuk penyimpanan data serta pengiriman notifikasi. Selain itu, monitoring telemetri memungkinkan pengawasan kondisi perangkat melalui parameter RSSI, uptime, free heap, reset reason, capture success rate, dan publish success rate sehingga mempermudah proses pemeliharaan dan identifikasi gangguan perangkat.
	
\end{enumerate}

\section{Saran}

Berdasarkan hasil penelitian yang telah dilakukan, terdapat beberapa saran yang dapat digunakan untuk pengembangan sistem pada penelitian selanjutnya.

\subsection{Pengembangan Sistem}

\begin{enumerate}
	
	\item Menambahkan fitur perekaman video berbasis event sehingga sistem tidak hanya menyimpan citra tetapi juga dapat merekam video ketika terdeteksi aktivitas pada area pengawasan.
	
	\item Menambahkan dukungan notifikasi melalui platform lain seperti Telegram, WhatsApp, atau email untuk meningkatkan fleksibilitas penyampaian informasi kepada pengguna.
	
	\item Mengembangkan dashboard monitoring yang lebih interaktif dengan menambahkan visualisasi statistik penggunaan perangkat, kondisi jaringan, dan histori aktivitas kamera.
	
	\item Menambahkan mekanisme manajemen pengguna yang lebih lengkap, seperti pembagian hak akses berdasarkan peran pengguna dan pengelolaan beberapa lokasi pengawasan dalam satu sistem.
	
\end{enumerate}

\subsection{Pengembangan Penelitian}

\begin{enumerate}
	
	\item Mengembangkan metode deteksi aktivitas yang lebih cerdas dengan memanfaatkan teknik pengolahan citra atau pembelajaran mesin untuk meningkatkan akurasi deteksi kejadian pada area pengawasan.
	
	\item Mengembangkan sistem analisis data telemetri untuk mendukung proses prediksi gangguan perangkat dan pemeliharaan preventif.
	
	\item Mengimplementasikan mekanisme edge computing yang lebih luas sehingga sebagian proses pengolahan data dapat dilakukan langsung pada perangkat atau node terdekat guna mengurangi beban server.
	
	\item Melakukan pengujian pada jumlah perangkat yang lebih besar untuk mengevaluasi skalabilitas sistem dalam lingkungan operasional yang lebih kompleks.
	
\end{enumerate}